% Generated from selected worked problems in committed QFT-soln sources.
% Source repository: https://github.com/Luca-Yucheng-Jin/QFT-soln
% Source commit: 80bd0650c1942eb3e25abaeacf7e934f20d9b8ef
\section{David Tong Statistical Field Theory, Sheet 2: Mostly Scaling}
\markboth{\thesection\quad TONG STATISTICAL FIELD THEORY, SHEET 2}{}

\noindent\textit{These prompts paraphrase David Tong's \emph{Statistical Field Theory},
Example Sheet 2 (November 2018). Problem numbers, starred status, equations,
and guidance follow the
\href{https://davidtong.org/pdfs/teaching/statistical-field-theory/99b.pdf}{official sheet}.}

\subsection{Sheet 2, Problem 1: Ornstein--Zernicke Correlations}

The Gaussian two-point function is
\begin{equation}
    G(x)=\int\frac{d^dk}{(2\pi)^d}
    \frac{e^{-ik\cdot x}}{\gamma k^2+\mu^2}
    =\frac{1}{\gamma}\int\frac{d^dk}{(2\pi)^d}
    \frac{e^{-ik\cdot x}}{k^2+\xi^{-2}},
    \qquad \xi^2=\frac{\gamma}{\mu^2}.
\end{equation}
Explain why it depends only on $r=|x|$. Establish the parameter integral
\begin{equation}
    \frac{1}{k^2+\xi^{-2}}
    =\int_0^\infty dt\,e^{-t(k^2+\xi^{-2})},
\end{equation}
perform the Gaussian momentum integral, and express the result, up to
an overall factor, as
\begin{equation}
    G(r)\sim\int_0^\infty dt\,e^{-S(t)},
    \qquad S(t)=\frac{r^2}{4t}+\frac{t}{\xi^2}+\frac d2\log t.
\end{equation}
Use a saddle-point approximation, stating the saddle equation and
quadratic prefactor, in the regimes $r\ll\xi$ and $r\gg\xi$.
Recover the Ornstein--Zernicke behaviours
\begin{equation}
    G(r)\sim r^{-(d-2)}\quad(r\ll\xi),
    \qquad
    G(r)\sim\frac{e^{-r/\xi}}{r^{(d-1)/2}}\quad(r\gg\xi).
\end{equation}
\begin{solution}
    Since the free energy respects global $SO(d)$ symmetry, the correlation function must also be invariant under rotation, and hence can only depend on $r$. The parameter integral is easy to see, and plugging into the correlation function, we find
    \begin{equation}
        \begin{aligned}
            G(r) &\sim\int d^dk \int_0^\infty dt \, e^{-i\mathbf{k}\cdot \mathbf{x} - t(k^2 +\xi^{-2})}\\
            &=\int d^dk \int_0^\infty dt \, e^{-(\sqrt{t}\mathbf{k} + i\mathbf{x}/2\sqrt{t})^2 - t\xi^{-2} - r^2/4t}\\
            &\sim\int^\infty_0dt \, \frac{e^{- t\xi^{-2} - r^2/4t}}{t^{d/2}}.
        \end{aligned}
    \end{equation}
    The saddle point is at
    \begin{equation}
        S'(t_*) = 0 \implies t_* = \frac{r^2}{d\pm \sqrt{d^2 + 4r^2/\xi^2}}.
    \end{equation}
    Since $t>0$, we take the plus sign. The correlation function becomes
    \begin{equation}
        G(r) \sim \int^\infty_0 dt \, e^{-S(t_*) - (t-t_*)^2S''(t_*)/2}\sim \frac{e^{-S(t_*)}}{\sqrt{S''(t_*)}}.
    \end{equation}
    For $r\ll \xi$, we have $t_* \sim r^2$, and hence
    \begin{equation}
        G(r) \sim r^{-(d-2)}.
    \end{equation}
    For $r \gg \xi$, we have $t_* \sim r\xi$, and hence
    \begin{equation}
        G(r) \sim \frac{e^{-r/\xi}}{r^{(d-1)/2}}.
    \end{equation}
\end{solution}

\subsection{Sheet 2, Problem 2: The Two-Point Function as a Green Function}

Verify that the correlation function satisfies
\begin{equation}
    (-\gamma\nabla^2+\mu^2)\langle\phi(x)\phi(0)\rangle
    \sim\delta(x),
\end{equation}
and explain its physical interpretation.
\begin{solution}
    This trivial to show. Physically, the correlation function tells us how does the system respond to a spiky source of perturbation.
\end{solution}

\subsection{Sheet 2, Problem 5*: Scaling in a Lifshitz Theory}

Separate $x=(x,\mathbf y)$ with $\mathbf y\in\mathbb R^{d-1}$, and
consider the free energy
\begin{equation}
    F[\phi]=\int d^dx\left[
    \frac12(\partial_x\phi)^2
    +\frac12(\nabla_{\!y}^{2}\phi)^2
    +\frac12\mu_0^2\phi^2\right].
\end{equation}
Write it in Fourier space using momenta $(k,\mathbf q)$. Under
$k'=\zeta k$, $\mathbf q'=\zeta^a\mathbf q$, and
$\phi'_{\mathbf{k}'}=\zeta^{-b}\phi_\mathbf{k}$, determine $a,b$ by normalising both
gradient terms, and find the running $\mu^2(\zeta)$.
In position space, obtain the Gaussian scaling dimension $\Delta_\phi$.
For an interaction $\int d^dx\,g_n\phi^{2n}$, find the dimension of
$g_n$ and show that $\phi^4$ is relevant for $d<7$ and irrelevant
for $d>7$.
\begin{solution}
    \begin{equation}
        F[\phi_\mathbf{k}] = \int \frac{d^dk}{(2\pi)^d}\frac{1}{2}(k^2 + \mathbf{q}^4 + \mu_0^2) \phi_\mathbf{k}\phi_{-\mathbf{k}} = \int \frac{d^dk'}{(2\pi)^d}\zeta^{-1-a(d-1)}\frac{1}{2}(\zeta^{-2}k'^2 + \zeta^{-4a}\mathbf{q}^4 + \mu^2_0)\zeta^{2b}\phi'_{-\mathbf{k}'}\phi_{\mathbf{k}'}'.
    \end{equation}
    In order to normalise both gradient terms, we must have
    \begin{equation}
        a = \frac{1}{2}, \quad b = \frac{d+5}{4}.
    \end{equation}
    Hence, we find
    \begin{equation}
        \mu^2(\zeta) = \zeta^2\mu_0^2.
    \end{equation}
    In position space, we have
    \begin{equation}
        \phi'(\mathbf{x}') = \int\frac{d^dk'}{(2\pi)^d}\phi_{\mathbf{k}'}'e^{ik'\cdot x'} = \zeta^{1+a(d-1)-b}\int \frac{d^dk}{(2\pi)^d}\phi_\mathbf{k}e^{    ik\cdot x} = \zeta^{\Delta_\phi}\phi(\mathbf{x})  \implies \Delta_\phi = \frac{d-3}{4}.
    \end{equation}
    For an interaction, we have
    \begin{equation}
        \int d^dx g_n \phi^{2n}(\mathbf{x}) = \int d^dx'\zeta^{1+a(d-1)}g_n\zeta^{-2n\Delta_\phi}\phi'^{2n}(\mathbf{x}') = \int d^dx'g_n(\zeta)\phi'^{2n}(\mathbf{x}'),
    \end{equation}
    which implies
    \begin{equation}
        g_n(\zeta) = \zeta^{-n(d-3)/2 +d/2+1/2}g_n.
    \end{equation}
    For $n=2$, we have
    \begin{equation}
        g_2(\zeta) = \zeta^{-(d-7)/2}g_2,
    \end{equation}
    and hence it is relevant for $d<7$.

\end{solution}

\subsection{Sheet 2, Problem 6: Scalar--Gauge Coupling}

For a complex scalar and gauge field, use
\begin{equation}
    F[\psi,A_i]=\int d^dx\left[
    \frac14F_{ij}F^{ij}
    +|\partial_i\psi-ieA_i\psi|^2+\mu^2|\psi|^2\right],
    \qquad F_{ij}=\partial_iA_j-\partial_jA_i.
\end{equation}
Find the critical dimension $d_c$ below which the coupling $e$ is
relevant and above which it is irrelevant. Discuss what may happen
at $d=d_c$. With a term $g|\psi|^4$, the $d=3$ case is the
Ginzburg--Landau theory of a superconductor.

\begin{solution}
    We have $[A_i] = [\psi] = -(d-2)/2$, and hence $[e] = (d-4)/2$ in length units. Clearly, the critical dimension is $d_c = 4$. Therefore, the coupling is relevant when $d<4$ and irrelevant when $d>4$. At $d=4$ it is marginal at this order; higher-order (loop) corrections determine whether it is marginally relevant or irrelevant.
\end{solution}

\subsection{Sheet 2, Problem 7: Emergent Rotational Symmetry}

A microscopic cubic lattice in $d$ dimensions has a local order
parameter $\phi(x)$ with $\phi\mapsto-\phi$ symmetry. Identify the
lowest-order interaction allowed by the lattice rotations but forbidden
by full $SO(d)$ rotations. Use scaling to explain why full rotational
invariance can emerge at long distances.
\begin{solution}
    The leading order contributing interaction is
    \begin{equation}
        \int d^dx \,\lambda\sum_{i=1}^d\phi \partial _i^4 \phi.
    \end{equation}
    The coupling constant now has a length dimension $[\lambda] = 2$, which is irrelevant for all $d$.
\end{solution}

\subsection{Sheet 2, Problem 8: Wick's Identity for a Gaussian Ensemble}

Let $G$ be a symmetric, positive-definite, invertible $n\times n$
matrix, and set
\begin{equation}
    \langle f(\phi)\rangle
    =\frac1{\mathcal N}\int_{\mathbb R^n}d^n\phi\,
    f(\phi)e^{-\phi\cdot G^{-1}\phi/2},
    \qquad \mathcal N=\det(2\pi G)^{1/2}.
\end{equation}
Prove $\langle\phi_a\phi_b\rangle=G_{ab}$ and, for constant $B_a$,
\begin{equation}
    \left\langle e^{B_a\phi_a}\right\rangle
    =\exp\!\left[\frac12B_a\langle\phi_a\phi_b\rangle B_b\right].
\end{equation}
Expand both sides to obtain all three pairings in
$\langle\phi_a\phi_b\phi_c\phi_d\rangle$. Show that odd moments
vanish and give the general even-moment pairing formula.
\begin{solution}
    Completing the square gives the generating function
    \begin{equation}
        Z(B)=\left\langle e^{B_a\phi_a}\right\rangle
        =e^{B_aG_{ab}B_b/2}.
    \end{equation}
    Differentiating twice at $B=0$ gives $\langle\phi_a\phi_b\rangle=G_{ab}$. Four derivatives give
    \begin{equation}
        \langle\phi_a\phi_b\phi_c\phi_d\rangle
        =G_{ab}G_{cd}+G_{ac}G_{bd}+G_{ad}G_{bc}.
    \end{equation}
    Since $Z(B)$ is even in $B$, all odd moments vanish. Repeated differentiation gives the sum over all pairings for every even moment. This is the finite-dimensional Gaussian version of the contractions in my \href{https://luca-yucheng-jin.github.io/luca-physics-site/notes/tong-qft-ps2.html#wicks-theorem-for-three-fields-q9}{Tong QFT Problem Sheet 2, Question 9 solution}.
\end{solution}
